% CLAUDE REPORT::WORLD-LEDGER-0001::2026-10-28T16:24Z
% !TEX program = xelatex
% Compiles with XeLaTeX or pdfLaTeX (ASCII source).
\documentclass[11pt]{article}
\usepackage[margin=1in]{geometry}
\usepackage{amsmath,amssymb,amsthm}
\usepackage{url}
\newtheorem{proposition}{Proposition}
\newtheorem{corollary}[proposition]{Corollary}
\newtheorem*{fact}{Fact}
\theoremstyle{remark}
\newtheorem*{remark}{Remark}
\newcommand{\Kt}{\mathrm{Kt}}
\newcommand{\Kpt}{K^{t}}
\newcommand{\R}{\mathcal{R}}
\newcommand{\cP}{\mathsf{P}}
\newcommand{\NP}{\mathsf{NP}}
\newcommand{\coNP}{\mathsf{coNP}}
\newcommand{\PSPACE}{\mathsf{PSPACE}}
\newcommand{\status}[1]{\textbf{[#1]}}

\title{CLAUDE WORLD-LEDGER-0001\\[2pt]
\large P vs NP world ledger; evaluation of the ``Pseudorandom Trap'' argument;\\
nested and time-budgeted SAT}
\author{CLAUDE (independent verification), from a chat session with the project owner}
\date{2026-10-28 09:24 America/Los\_Angeles}

\begin{document}
\maketitle

\noindent\textbf{Status.} P versus NP remains OPEN. Nothing in this report proves or claims a separation.
This is an additive reference resource requested by the project owner. It changes no task schedules,
agent count, role ownership or scientific allocation, and it does not edit root \texttt{main.tex}.
It was not produced as an LSPV run and did not perform the LSPV corpus-wide reconciliation.

\noindent\textbf{Status labels} follow \texttt{AGENTS.md}: PROVED/EXACT, KNOWN THEOREM, COMPUTATIONAL,
CONDITIONAL, CANDIDATE, OPEN, FALSIFIED/RETIRED.

\section{The artifact}

\noindent Location: \texttt{research/agents/claude/world-ledger/}.

\begin{itemize}
  \item \texttt{build\_ledger.py} --- source of truth (the \texttt{rows} list); regenerates both outputs.
  \item \texttt{P\_vs\_NP\_World\_Ledger.xlsx} --- spreadsheet with live Category and Consistency formulas.
  \item \texttt{world\_ledger.csv} --- plain-text copy; Category and Consistency computed in Python by
        the same rules.
  \item \texttt{README.md} --- how to read it and how to add rows.
\end{itemize}

Each row is a statement $S$ marked True, False or Open in Impagliazzo's five worlds:
Algorithmica ($\cP=\NP$); Heuristica ($\cP\neq\NP$, NP easy on average);
Pessiland (NP hard on average, no one-way functions); Minicrypt (one-way functions, no public-key
encryption); Cryptomania (public-key encryption). Each world cell is a CONDITIONAL claim ``if world $W$
then $S$''. The five values determine the category:

\begin{center}
\begin{tabular}{lll}
Algorithmica & four $\cP\neq\NP$ worlds & Category \\ \hline
True & all True & Theorem (win-win) \\
True & all False & Equivalent to $\cP=\NP$ \\
True & no False, some Open & Half-win (true if $\cP=\NP$) \\
True & some (not all) False & World separator (true if $\cP=\NP$) \\
False & all False & Refuted in every world \\
False & all True & Equivalent to $\cP\neq\NP$ \\
False & both True and False & World separator (implies $\cP\neq\NP$) \\
False & otherwise & Implies $\cP\neq\NP$ (open) \\
Open & all False & Implies $\cP=\NP$ (open) \\
Open & all True & Half-win (true if $\cP\neq\NP$) \\
Open & both True and False & World separator (open if $\cP=\NP$) \\
Open & otherwise & Undetermined by P vs NP \\
\end{tabular}
\end{center}

A Consistency column flags CONFLICT whenever a world contradicts a row whose known status is Theorem
(or Refuted). A genuine conflict would eliminate that world; in practice it signals an entry error.

\status{COMPUTATIONAL} Verification of the artifact: LibreOffice recalculation of all 70 formulas with
0 errors; the Python classifier and the recalculated spreadsheet agree on 28/28 rows; setting
Algorithmica to False on a Theorem row fires the CONFLICT flag. Caveat: the five worlds are an
exhaustive split only under fixed definitions; infinitely-often and non-uniform variants can fall
between them.

\section{Definitions used by the cited results}

Fix a universal machine $U$ with polynomial simulation overhead and a constant $c\ge 0$.
$K(x)=\min\{|p|:U(p)=x\}$; for a polynomial $t$, $\Kpt(x)=\min\{|p|:U(p)\text{ outputs }x\text{ within }t(|x|)\text{ steps}\}$;
$\Kt(x)=\min\{|p|+\lceil\log_2 T\rceil: U(p)\text{ outputs }x\text{ within }T\text{ steps}\}$;
conditional $\Kt(x\mid A)$ gives $U$ the input $A$.
For $\mu\in\{K,\Kt,\Kpt\}$ let $\R(\mu)=\{x:\mu(x)\ge |x|-c\}$.
A set is \emph{P-immune} if it is infinite and has no infinite subset in $\cP$; \emph{immune} if it has
no infinite c.e.\ subset.

\begin{fact}[counting] \status{PROVED/EXACT}
For $n\ge c$, at least $2^n-2^{n-c}+1$ strings of length $n$ lie in $\R(K)$, since only $2^{n-c}-1$
programs are shorter than $n-c$. As $K\le \Kt$ and $K\le\Kpt$, all three sets $\R(\mu)$ are infinite.
\end{fact}

\section{Evaluation of the ``Pseudorandom Trap'' argument (ledger rows L20--L24)}

The argument (submitted by the owner, 2026-10-28) claims: (A) $\cP=\NP$ yields an infinite
polynomial-time subset of the random strings; (B) the random strings are unconditionally P-immune;
hence $\cP\neq\NP$. The inference is valid; the verdict concerns the lemmas.

\begin{proposition}[row L22] \status{KNOWN THEOREM}
$\R(K)$ is immune. Hence Lemma A for $K$ has an impossible conclusion and is equivalent to $\cP\neq\NP$.
\end{proposition}
\begin{proof}[Sketch]
If $W\subseteq\R(K)$ is infinite and c.e., let $y_m$ be the first enumerated element of length $\ge m$.
Then $m-c\le K(y_m)\le \log_2(m+1)+O(1)$, false for large $m$.
\end{proof}

\begin{proposition} \status{PROVED/EXACT}
If $\cP=\NP$ then $\R(\Kt)$ is P-immune; in particular $\R(\Kt)\notin\cP$. Hence Lemma A for $\Kt$ is
equivalent to $\cP\neq\NP$.
\end{proposition}
\begin{proof}
Let $L\in\cP$ be infinite. Under $\cP=\NP$, prefix search (each test an NP query) finds the
lexicographically first $x_n\in L$ of length $n$ in time $n^k$ whenever one exists. Then
$\Kt(x_n)\le (k+1)\log_2 n+O(1)<n-c$ for infinitely many $n$, so $L\not\subseteq\R(\Kt)$.
\end{proof}

\begin{proposition}[rows L20, L21] \status{PROVED/EXACT}
For every polynomial $t$: (i) $\R(\Kpt)\in\coNP$ and is infinite; (ii) if $\cP=\NP$ then
$\R(\Kpt)\in\cP$, so Lemma A for $\Kpt$ holds trivially; (iii) if $\R(\Kpt)$ is P-immune then
$\cP\neq\NP$.
\end{proposition}
\begin{proof}
(i) A program $p$ with $|p|<|x|-c$ that prints $x$ within $t(|x|)$ steps is an NP witness for
$x\notin\R(\Kpt)$. (ii) $\cP=\NP$ gives $\coNP=\cP$. (iii) P-immunity gives $\R(\Kpt)\notin\cP$,
so $\coNP\neq\cP$.
\end{proof}

\noindent\textbf{Verdict.} \status{FALSIFIED/RETIRED as a proof.} In every reading, one lemma is
equivalent to or implies $\cP\neq\NP$. Lemma B for polynomial $t$ (``$\R(\Kpt)$ is P-immune'') is
\status{OPEN} and implies $\cP\neq\NP$; no unconditional theorem of this form is known, and any such
theorem would itself prove the separation. Relativization check: for a PSPACE-complete oracle $A$,
$\cP^A=\NP^A$ and $\R(\Kpt)$ relative to $A$ lies in $\coNP^A=\cP^A$, so relativized Lemma B is false
while every other step relativizes.

\begin{proposition}[row L20] \status{PROVED/EXACT}
If $\R(\Kpt)\in\cP$ for every polynomial $t$, then no PRG $G:\{0,1\}^n\to\{0,1\}^{2n}$ computable in
polynomial time is secure; with HILL (one-way functions give PRGs, \status{KNOWN THEOREM}) one-way
functions do not exist.
\end{proposition}
\begin{proof}
For $t$ large enough, every output satisfies $\Kpt(G(s))\le n+O(1)<2n-c$, so outputs never lie in
$\R(\Kpt)$, while a uniform $2n$-bit string lies in it with probability at least $1-2^{-c}$.
\end{proof}

\begin{proposition}[Chaitin; row L24] \status{KNOWN THEOREM}
For a consistent, computably axiomatized theory $T$ proving all true $\Sigma_1$ sentences, there is
$L_T$ such that $T$ proves ``$K(x)>L_T$'' for no $x$. It limits proof, not existence (row L23 is
refuted by the counting fact).
\end{proposition}

\section{Nested SAT (rows L09--L12, L27)}

Setting: the satisfiability answer of an instance $A$ appears as a literal $a$ in clauses of $B$.

\begin{enumerate}
\item \status{PROVED/EXACT} If inner satisfying assignments feed outer clauses, the whole is one SAT
instance: NP.
\item \status{PROVED/EXACT} \emph{Positive occurrences only} (any polynomial DAG of nested instances): NP.
Certificate: a claimed bit $b_i$ for every nested instance, a witness for each $b_i=1$ (with inner bits
replaced by their claimed values), and a witness for the outermost formula. Soundness: in topological
order $b_i\le a_i$, because each formula is monotone in its nested bits. Completeness: claim the true bits.
For one level, $B[a]$ is satisfiable iff $B[0]$ is, or $A$ and $B[1]$ are.
\item \status{KNOWN THEOREM / PROVED/EXACT (sketch)} \emph{Negated occurrences, fixed instances}: $d$
levels lie in $\cP^{\NP[d]}$ ($d$ adaptive queries); polynomially many levels are $\cP^{\NP}$-complete
(encode the polynomial-time query generator of a $\cP^{\NP}$ machine by Cook--Levin, with earlier answers
as literals).
\item \status{KNOWN THEOREM} \emph{Inner formulas depend on outer variables}: depth $k$ gives
$\Sigma_k^{\mathrm p}$-completeness; polynomial depth gives TQBF, $\PSPACE$-complete.
\item \status{OPEN} $\cP=\mathsf{PH}$ is equivalent to $\cP=\NP$ and collapses every fixed depth. It does
not obviously collapse polynomial depth: substituting a size-$n^{c}$ encoding of an inner solver at each
level gives size about $n^{c^d}$ after $d$ levels. $\cP=\NP\neq\PSPACE$ is not refuted by any known
result. (Remark only; not a proof that $\cP=\NP$ fails to imply $\cP=\PSPACE$.)
\end{enumerate}

\section{Time-budgeted SAT and the Kt equivalence (row L28)}

Proposal: $a$ = ``algorithm $M$ finds a satisfying assignment of $A$ within $|A|^2$ steps'', used as a
literal in $B$.

\begin{enumerate}
\item \status{PROVED/EXACT} For a fixed deterministic $M$, $a$ is polynomial-time computable, so the
nested problem is in NP (NP-complete, as $B$ alone is 3-SAT). $a=1$ is certified; $a=0$ only means
``not found in time''.
\item \status{PROVED/EXACT} A nondeterministic budget is vacuous: a guessed assignment is checked in
linear time, so $a=\mathrm{SAT}(A)$ and Section 4 applies.
\item \status{PROVED/EXACT} If $\cP\neq\NP$, every fixed polynomial-time $M$ errs on infinitely many
satisfiable inputs (finitely many errors could be patched by table lookup). \status{KNOWN THEOREM}
Gutfreund--Shaltiel--Ta-Shma (2007): under $\cP\neq\NP$, failing instances for any such deterministic
$M$ can be generated efficiently.
\item Literature status (van Melkebeek survey): no superlinear deterministic time lower bound for SAT
on general random-access machines is known without a space restriction, so ``SAT is decidable in
$n^2$ steps'' is not refuted. Ledger classification of that statement: Implies $\cP=\NP$ (open).
\end{enumerate}

\begin{proposition}[row L28] \status{PROVED/EXACT}
$\cP=\NP$ if and only if there is a constant $c$ such that every satisfiable formula $A$ has a
satisfying assignment $x$ with $\Kt(x\mid A)\le c\log_2|A|$.
\end{proposition}
\begin{proof}
($\Rightarrow$) By self-reducibility, $\cP=\NP$ gives a search algorithm $S$ running in time
$|A|^k$; its output $x$ has $\Kt(x\mid A)\le |S|+k\log_2|A|$.
($\Leftarrow$) Universal search on input $A$: in phase $i$ run every program $p$ with $|p|\le i$ for
$2^{\,i-|p|}$ steps and check each output. A solution with $\Kt(x\mid A)\le c\log_2 n$ is found by
phase $c\log_2 n+O(1)$, at total cost $2^{O(c\log n)}\mathrm{poly}(n)$. Reject if nothing is found by
then. This decides SAT in polynomial time.
\end{proof}

\section{Relevance to the current ARC/FLMV allocation (row L13)}

\status{KNOWN THEOREM} Williams: SAT $\notin\mathrm{DTISP}(n^{c},n^{o(1)})$ for $c<2\cos(\pi/7)\approx 1.8019$;
Buss--Williams (2015): alternation-trading proofs built from the known speedup/slowdown rules cannot
reach $c\ge 2\cos(\pi/7)$, and they compute the achievable time-space tradeoffs for $n^{\varepsilon}$
space, $0<\varepsilon<1$.
\status{CANDIDATE} relevance note: agents pursuing the strict FLMV box should check whether their route
is an alternation-trading proof in the Buss--Williams sense; if it is, the Buss--Williams limits for
$0<\varepsilon<1$ apply to it, and escaping them needs a rule outside the two standard schemas.
This report does not assess whether the current box lies inside or outside those limits.

\section{Literature}

Verified against primary pages in this session:
\begin{itemize}
\item L.~Chen, S.~Hirahara, H.~Ren, Symmetric exponential time requires near-maximum circuit size,
      \url{https://arxiv.org/abs/2309.12912}.
\item Z.~Li, Symmetric exponential time requires near-maximum circuit size: simplified, truly uniform,
      \url{https://arxiv.org/abs/2310.17762}.
\item S.~Buss, R.~Williams, Limits on alternation-trading proofs for time-space lower bounds,
      Computational Complexity 24(3), 2015; \url{https://eccc.weizmann.ac.il/report/2011/031}.
\item D.~Gutfreund, R.~Shaltiel, A.~Ta-Shma, If NP languages are hard on the worst-case, then it is easy
      to find their hard instances, Computational Complexity, 2007;
      \url{https://link.springer.com/article/10.1007/s00037-007-0235-8}.
\item D.~van Melkebeek, A survey of lower bounds for satisfiability and related problems,
      \url{https://pages.cs.wisc.edu/~dieter/Research/sat-lb-survey.html}.
\end{itemize}
Cited from memory (verify before load-bearing use): Impagliazzo 1995 (five worlds); Razborov--Rudich 1997;
Baker--Gill--Solovay 1975; Chaitin 1974; Liu--Pass 2020; Levin 1973; Kannan 1982; Karp--Lipton 1980;
Ladner 1975; Mahaney 1982; Kabanets--Cai 2000; H{\aa}stad--Impagliazzo--Levin--Luby 1999;
Impagliazzo--Wigderson 1997; Sipser 1983 and Lautemann 1983; Fortnow--Lipton--van Melkebeek--Viglas 2005;
Williams 2011 (NEXP $\not\subseteq$ ACC$^0$).

\section{Open items handed forward}

\begin{itemize}
\item Technique tags marked Unverified in the ledger (L13, L14, L19, L26) need checking against the
      primary papers.
\item Lemma B for polynomial $t$ (row L21) is OPEN and implies $\cP\neq\NP$; it is not a usable premise.
\item Half-wins L03--L06: the productive targets are those whose open side would need a
      non-relativizing step.
\end{itemize}

\end{document}
% END_OF_VERIFIED_REPORT::CLAUDE::WORLD-LEDGER-0001
